\section{迭代简化 Harness}

第一版 harness 效果令人鼓舞，但也笨重、缓慢、昂贵。下一步的逻辑是：在不降低性能的前提下简化 harness。

\begin{importantbox}{Harness 简化的核心原则}
\textit{``Find the simplest solution possible, and only increase complexity when needed.''} ——来自 Anthropic 的 Building Effective Agents 博客文章

实际操作上，作者首先尝试了激进的简化但失败了（无法复现原有性能），随后转向更系统的方法：\textbf{每次只移除一个组件}，观察其对最终结果的影响。
\end{importantbox}

\subsection{Opus 4.6 带来的变化}

在迭代过程中，Opus 4.6 发布，其关键改进直接对应了 harness 之前需要补偿的能力缺陷：

\begin{itemize}
    \item 更细致的规划能力
    \item 更持久的 agentic task 执行
    \item 更可靠的大型代码库操作
    \item 更好的代码审查和调试（自我纠错）能力
    \item 显著改善的 long-context retrieval
\end{itemize}

这些改进意味着之前 harness 中的许多组件可能不再是必要的。

\subsection{移除 Sprint 结构}

作者首先移除了 sprint 构造。Sprint 结构的原始目的是将工作分解为模型可以连贯处理的小块。Opus 4.6 的能力提升使得模型可能无需这种分解就能原生处理任务。

关键变化：
\begin{itemize}
    \item \textbf{Planner 保留}：没有 planner，generator 会不充分地确定范围，直接开始构建，产出功能较少的应用
    \item \textbf{Evaluator 移至末尾}：从每个 sprint 评分改为整体运行结束后的单次评估
    \item \textbf{Evaluator 的价值变为任务依赖}：在 4.5 上，构建处于 generator 能力边缘，evaluator 普遍有价值；在 4.6 上，模型能力边界外移，许多任务 generator 已能独立处理好，evaluator 仅在超出 generator 能力边界的任务部分仍有提升
\end{itemize}

\begin{warningbox}{Evaluator 不是固定的``有/无''决策}
Evaluator 的价值取决于任务是否处于当前模型能力的边缘。模型越强，evaluator 有价值的任务范围越小，但在那些仍处于边缘的任务上，evaluator 依然提供显著提升。这意味着随着模型演进，harness 设计者需要持续重新评估每个组件的贡献。
\end{warningbox}

\subsection{DAW 构建实验}

使用更新后的 harness，作者测试了一个更具挑战性的 prompt：

\begin{quote}
\textit{``Build a fully featured DAW in the browser using the Web Audio API.''}
\end{quote}

\begin{table}[H]
\centering
\begin{tabular}{lrr}
\toprule
\textbf{Agent \& 阶段} & \textbf{时长} & \textbf{成本} \\
\midrule
Planner & 4.7 min & \$0.46 \\
Build (Round 1) & 2 hr 7 min & \$71.08 \\
QA (Round 1) & 8.8 min & \$3.24 \\
Build (Round 2) & 1 hr 2 min & \$36.89 \\
QA (Round 2) & 6.8 min & \$3.09 \\
Build (Round 3) & 10.9 min & \$5.88 \\
QA (Round 3) & 9.6 min & \$4.06 \\
\midrule
\textbf{Total} & \textbf{3 hr 50 min} & \textbf{\$124.70} \\
\bottomrule
\end{tabular}
\caption{DAW 构建各阶段的时长与成本（Updated Harness + Opus 4.6）}
\end{table}

关键观察：

\begin{itemize}
    \item Builder 在没有 sprint 分解的情况下连贯运行超过 \textbf{2 小时}——这在 Opus 4.5 上是不可能的
    \item QA 仍然发现了真实缺陷，例如第一轮反馈指出：多个核心 DAW 功能仅为展示性而无交互深度（clip 不能在 timeline 上拖动，无乐器 UI 面板，无视觉效果编辑器）
    \item 第二轮 QA 继续发现：音频录制仍为 stub、clip 缩放和分割未实现、效果可视化仅为数字滑块而非图形化
    \item 最终产出是一个功能完整的浏览器 DAW：工作的 arrangement view、mixer 和 transport，用户可以完全通过 prompt 与内置 AI agent 交互来创作歌曲
\end{itemize}

\begin{table}[H]
\centering
\begin{tabular}{lrrp{5cm}}
\toprule
\textbf{Harness 版本} & \textbf{时长} & \textbf{成本} & \textbf{特点} \\
\midrule
Solo agent (4.5) & 20 min & \$9 & 核心功能不可用 \\
Full harness (4.5) & 6 hr & \$200 & 核心功能可用，有粗糙之处 \\
Updated harness (4.6) & 3 hr 50 min & \$125 & 更简洁，同等或更高质量 \\
\bottomrule
\end{tabular}
\caption{三个版本 harness 的综合对比}
\end{table}

\subsection{本章小结}

从 Opus 4.5 到 4.6 的模型升级使得 harness 可以显著简化：sprint 结构被移除，evaluator 改为末尾单次评估。成本从 \$200/6hr 降至 \$125/4hr，同时保持甚至提升了产出质量。核心方法论是：系统地逐个移除组件并观察影响，而非激进地一次性简化。

